%\newpage
\section{The Hardness of {\sc oracle exact matroid basis}}
\label{sec:technical}
\label{sec:Pi}
In this section, we prove \Cref{thm:emb}. We use in the proof the family of $\Pi$-{\em matroids}. A member in the $\Pi$-matroid family is given by four arguments: $n,k,\alpha \in \mathbb{N}_{>0}$, and $\Pi \subseteq 2^{[n]}$. For any~$m\in \mathbb{N}$, let $[m]=\{1,\ldots, m\}$.
 The first argument, $n \in \mathbb{N}_{>0}$, is the number of elements, and the ground set is $[n]$. The second argument, $k \in [n]$, is the rank of the matroid. 
% \footnote{The rank of a matroid is the maximum cardinality of an independent set of the matroid.} 
 The third argument, $\alpha\in \mathbb{N}_{>0}$, is a target value, that is usually equal to $\textnormal{\textsf{sum}}(S)$ for some $S \subseteq [n]$, where $\textnormal{\textsf{sum}}(S) = \sum_{i \in S} i$. The last argument is a family of subsets $\Pi \subseteq 2^{[n]}$, where $2^{[n]} = 2^{\{1,\ldots,n\}}$.


The set $\Pi$ is called the {\em secret family} because finding $S\in \Pi$ is possible only via repeated queries to the membership oracle of the matroid. Since $\Pi$ can have an arbitrary structure, this may require exhaustive enumeration.


\begin{definition}
	\label{def:Matroid}
	Let $n,k,\alpha \in \mathbb{N}_{>0}$. For some $\Pi \subseteq 2^{[n]}$, define the {\bf $\Pi$-matroid} on $n,k$, and $\alpha$ as $M_{n,k,\alpha}(\Pi) = ([n],\cI_{n,k,\alpha}(\Pi))$, where $$\cI_{n,k,\alpha}(\Pi) = \mathcal{J}_{n,k} \cup \mathcal{K}_{n,k,\alpha} \cup \mathcal{L}_{n,k,\alpha}(\Pi)$$ and $\mathcal{J}_{n,k}, \mathcal{K}_{n,k,\alpha}, \mathcal{L}_{n,k,\alpha}(\Pi)$ are defined as follows. \begin{equation}
		\label{eq:1}
		\begin{aligned}
			\mathcal{J}_{n,k}~~~~~~ = {} & \big\{S \subseteq [n]~\big|~ |S| < k\big\} \\
			\mathcal{K}_{n,k,\alpha}~~~~ = {} & \big\{S \subseteq [n]~\big|~ |S| = k, \textnormal{\textsf{sum}}(S) \neq \alpha\big\} \\
			\mathcal{L}_{n,k,\alpha}(\Pi) = {} & \big\{S \subseteq [n]~\big|~ |S| = k, \textnormal{\textsf{sum}}(S) = \alpha, S \in \Pi\big\}. \\
		\end{aligned}
	\end{equation}
\end{definition} 

\begin{figure}
	\centering
	\vspace{-4cm} 
	\begin{tikzpicture}[scale=1.4, every node/.style={draw, circle, inner sep=1pt}]
		\node (11) at (5.5,4) {$\bf \textcolor{teal}{1}$};
		\node (22) at (7,4) {$\bf \textcolor{black}{2}$};
		\node (33) at (5.5,3) {$\bf \textcolor{red}{3}$};
		\node (44) at (7,3) {$\bf \textcolor{blue}{4}$};
		\draw (11) -- (22);
		\draw (11) -- (33);
		\draw (11) -- (44);
		\draw (22) -- (44);
		\draw (33) -- (44);
		
		\node[draw=none] at (9, 3.5) {$~\Pi = \{\textnormal{independent sets of } G\}$};
		
		\node[draw=none] at (4.5, 3.5) {A graph $G$};
		% Formula
		\node[draw=none] at (6.75, 2.5) {$\mathcal{J}_{4,2} = \{\emptyset, \{$\bf \textcolor{teal}{1}$\},\{$\bf \textcolor{black}{2}$\},\{$\bf \textcolor{red}{3}$\},\{$\bf \textcolor{blue}{4}$\}\}, ~~\mathcal{K}_{4,2,5} = \{\{$\bf \textcolor{black}{2}$,$\bf \textcolor{blue}{4}$\},\{$\bf \textcolor{black}{2}$,$\bf \textcolor{teal}{1}$\},\{$\bf \textcolor{blue}{4}$,$\bf \textcolor{red}{3}$\},\{$\bf \textcolor{red}{3}$,$\bf \textcolor{teal}{1}$\}\}, ~~\mathcal{L}_{4,2,5}(\Pi) = \{\{$\bf \textcolor{black}{2}$,$\bf \textcolor{red}{3}$\}\}$};
	\end{tikzpicture}
	\vspace{-7.5cm} 
	\caption{\label{fig:Pi} The independent sets of the $\Pi$-matroid $M_{n,k,\alpha}(\Pi)$, with parameters $n = 4$, $k = 2$, and $\alpha = 5$. The secret family $\Pi$ contains all independent sets in the graph $G$, where $\{2,3\}$ is the only independent set in $G$ with $k$ elements.}
	
\end{figure}


In words, $\mathcal{J}_{n,k}$ contains all subsets of strictly less than $k$ elements;  $\mathcal{K}_{n,k,\alpha}$ contains all subsets of cardinality~$k$ whose total sum is not~$\alpha$; finally, 	$\mathcal{L}_{n,k,\alpha}(\Pi)$ contains all subsets of cardinality~$k$ and total sum~$\alpha$ which also belong to $\Pi$. See \Cref{fig:Pi} for an example of a member of the $\Pi$-matroid family. Using a simple argument, we show that the set system in 
\Cref{def:Matroid} is indeed a matroid.
For the sets 
$\cI_{n,k,\alpha}(\Pi)$, $ \mathcal{J}_{n,k}$, $\mathcal{K}_{n,k,\alpha}$ and $\mathcal{L}_{n,k,\alpha}(\Pi)$ defined in \Cref{def:Matroid},
we often omit the subscripts ${n,k,\alpha}$ and ${n,k}$ when the values of $n,k,\alpha$ are known by context. For simplicity, for any set $A$ and an element $a$, let $A+a$, $A-a$ be $A \cup \{a\}$ and $A \setminus \{a\}$, respectively. 




\begin{restatable}{lemma}{isamatroid}
	\label{lem:1}
	For every  $n,k,\alpha \in \mathbb{N}_{>0}$, and $\Pi \subseteq 2^{[n]}$ it holds that $\mm$ is a matroid.
\end{restatable}
	\begin{proof}
	We first note that $\emptyset \in \mathcal{J}$ since $0<k$; therefore, $\emptyset \in \cI(\Pi)$. For the hereditary property, let $A \in \cI(\Pi)$. For all $B \subset A$ it holds that $|B| < k$; thus, $B \in \mathcal{J}$ and it follows that $B \in \cI(\Pi)$. For the exchange property, let $A,B \in \cI(\Pi)$ such that $|A| > |B|$. We consider the following cases. 
	\begin{enumerate}
		\item $|B| < k-1$. Then, for all $e \in A \setminus B$ it holds that $|B+e| < k-1+1 = k$; thus, $B+e \in \cJ$ and it follows that $B+e \in \cI(\Pi)$. Note that there is such $e  \in A \setminus B$ because $|A| > |B|$. 
		
		\item $|B| = k-1$ and $|A| = k$. We consider two subcases.  
		\begin{enumerate}
			\item $B \subseteq A$. Then, as $|A| > |B|$ there is $e \in A \setminus B$; thus, $B+e = A$ (because $|B| = k-1$ and $|A| = k$). As $A \in \cI(\Pi)$, it follows that $B+e \in \cI(\Pi)$.
			
			\item $B \not\subseteq A$.  Then, as $|B| = k-1$ and $|A| = k$ it follows that $|B \cap A| < |B| = k-1$; thus, $|A \setminus B| = |A|-|A \cap B| > k - (k-1) = 1$. Therefore, there are $e,f \in A \setminus B$ such that $e \neq f$. % thus, %since $f$ is a bijection, 
			%	it holds that $f(u) \neq f(v)$. 
			It follows that there is $g \in \{e,f\}$ such that $\textnormal{\textsf{sum}}(B+g) = \textnormal{\textsf{sum}}(B)+g \neq \alpha$. We conclude from \eqref{eq:1} that $B+g \in \mathcal{K}$,
			implying that $B+g \in \cI(\Pi)$.
		\end{enumerate}
\end{enumerate}	\end{proof}
%The proof of \Cref{lem:1} is given in \Cref{sec:Pi}. 
Observe that $\mathcal{K}_{n,k,\alpha} \cup \mathcal{L}_{n,k,\alpha}(\Pi)$ is the set of bases of the matroid. Moreover, for any arguments $n$, $k$, $\alpha$ and $\Pi$,
the cardinality of every  dependent set $S\in 2^{[n]}\setminus \cI_{n,k,\alpha}(\Pi)$  is at least the rank of the $\Pi$-matroid $M_{n,k,\alpha} (\Pi)=([n],\cI_{n,k,\alpha}(\Pi))$.
Such matroids are known as {\em paving matroids} (see, e.g.,~ \cite{Ox2006,PvdP14}).
%In this section we give the hardness result for EMB as stated in \Cref{thm:emb}. For the proof of \Cref{thm:emb}, 
Using $\Pi$-matroids, we define the following collection of oracle-EMB instances. In these instances, the matroid is a $\Pi$-matroid, 
%. goal is to decide if $\mathcal{L}_{n,k,\alpha}(\Pi) \neq \emptyset$, 
where $\Pi$ is some unknown fixed family of subsets of the ground set. 
\begin{definition}
	\label{def:EMBpi}
	For every $n,k,\alpha \in \mathbb{N}_{>0}$, $\Pi \subseteq 2^{[n]}$ define the 
	{\em induced} oracle-\textnormal{EMB} instance of $n,k,\alpha,\Pi$, denoted $I_{n,k,\alpha}(\Pi)$, as follows. Let
	$\textnormal{\textsf{id}}_n: [n] \rightarrow [n]$, where $\textnormal{\textsf{id}}_n(i) = i~\forall i \in [n]$. Then, 
	$I_{n,k,\alpha}(\Pi) = ([n],\cI_{n,k,\alpha}(\Pi),\textnormal{\textsf{id}}_n,\alpha)$.
\end{definition}

Observe that the above is indeed an oracle-EMB instance if the independent sets of the given matroid are accessible via a membership oracle. The following is an easy consequence of \Cref{def:EMBpi}.

\begin{observation}
	\label{obs:YesNo}
	For every $n,k,\alpha \in \mathbb{N}_{>0}$ and $\Pi \subseteq 2^{[n]}$, it holds that $I_{n,k,\alpha}(\Pi)$ is an \textnormal{oracle-EMB} ''yes"-instance if and only if there is $S \in \Pi$ such that $|S| = k$ and $\textnormal{\textsf{sum}}(S) = \textnormal{\textsf{id}}_n(S) = \alpha$. 
\end{observation}

By \Cref{obs:YesNo}, an algorithm that finds an independent set of $\mathcal{M}_{n,k,\alpha}(\Pi)$ satisfying $|S| = k$ and $\textnormal{\textsf{sum}}(S) = \alpha$, %(i.e., deciding the EMB instance $I_{n,k,\alpha}(\Pi)$), 
in fact outputs a subset $S \in \Pi$. 
As the input for an induced oracle-EMB instance $I_{n,k,\alpha}(\Pi)$ does not contain an explicit encoding of $\Pi$, finding $S\in \Pi$ requires a sequence of queries to the membership oracle of $\mm$.  Roughly speaking, to decide $I_{n,k,\alpha}(\Pi)$, 
an algorithm for oracle-EMB must iterate over all (exponentially many) subsets $S \subseteq [n]$ such that $|S| = k$ and $\textnormal{\textsf{sum}}(S) = \alpha$. This is the intuition
behind the proof of the next result.  
